\documentclass[11pt]{article}

\usepackage[T1]{fontenc}
\usepackage{lmodern}
\usepackage{amsmath}
\usepackage{booktabs}
\usepackage{natbib}

\title{{{title}}}
\author{{{author}}}
\date{\today}

\begin{document}

\maketitle

\begin{abstract}
State the question, what you did, and what you found, in two or three sentences.
\end{abstract}

\section{Introduction}
\label{sec:intro}

Start writing here. Cite a source in the sentence, as \citet{knuth1984literate}
do, or in parentheses \citep{lamport1994latex}. Section~\ref{sec:method} shows an
equation and a table.

\section{Method}
\label{sec:method}

Equation~\eqref{eq:mean} defines the mean of $n$ observations:
\begin{equation}
  \bar{x} = \frac{1}{n} \sum_{i=1}^{n} x_i .
  \label{eq:mean}
\end{equation}
Table~\ref{tab:symbols} lists the symbols.

\begin{table}[ht]
  \centering
  \caption{Symbols used in Equation~\eqref{eq:mean}.}
  \label{tab:symbols}
  \begin{tabular}{cl}
    \toprule
    Symbol    & Meaning \\
    \midrule
    $n$       & number of observations \\
    $x_i$     & the $i$-th observation \\
    $\bar{x}$ & their mean \\
    \bottomrule
  \end{tabular}
\end{table}

\bibliographystyle{plainnat}
\bibliography{refs}

\end{document}
